\begin{afterword}
\summaryhead

The post defines the availability heuristic: judging how frequent or likely an event is by how easily examples come to mind. Its evidence is a 1978 study in which people judged accidents to kill about as many people as disease, and homicide to be more common than suicide. A 1979 follow-up found that these errors matched what two newspapers reported.

The author argues that selective reporting is a major source of such errors, and a stronger one now than for our ancestors, who learned most things first- or second-hand. Bill Gates is heard of far more often than his numbers justify, and the very poor far less.

Availability also seems to make people treat events they have never experienced as impossible. People refuse flood insurance even when it is heavily subsidized, and dams, by preventing small floods, lead to fewer precautions and greater average yearly damage. Past small hazards set the perceived limit on risk. Memory is not always a good guide to probability.

\responsehead

The post is a short summary of a real research finding, and it reports its central example correctly. In the 1978 study about 70 percent of subjects did pick homicide as the more frequent cause of death, and accidents were judged about as deadly as disease. The problems are in the claims the post adds.

Two passages add claims that nothing in the post supports. The account of the ancestral environment, with ``at most one layer of selective reporting,'' is a claim with no source. The Bill Gates example is about whom people compare themselves with, not about the judgment of frequency the post defined.

The flood section misattributes a quotation and overstates a finding. Much of the post follows the author's own paper on global risks. The draft of that paper, like the post, gives the flood quotation to Kunreuther and colleagues; the version published in 2008 gives it to Robert Kates, writing in 1962. The post gives its readers the earlier attribution. Then the text moves from ``reduced precautions,'' the finding it attributes to the cited book, to the rule that a protected society ``takes no action against major risks.'' The post turns a finding about one kind of hazard into a law about societies. The last line returns to the modest claim that memory is ``not always'' a good guide, which the evidence does support.

In short: a fair summary of a real finding, with a quotation whose source the author's own later paper changes, and a rule about societies drawn from floods.
\end{afterword}
